\documentclass[11pt,reqno]{amsart}

\usepackage[T1]{fontenc}
\usepackage{lmodern}
\usepackage{microtype}
\usepackage{mathtools}
\usepackage{amssymb}
\usepackage{xcolor}
\usepackage[colorlinks=true,linkcolor=blue!55!black,citecolor=blue!55!black,urlcolor=blue!55!black]{hyperref}
\usepackage[nameinlink,capitalise,noabbrev]{cleveref}

\newtheorem{theorem}{Theorem}[section]
\newtheorem{proposition}[theorem]{Proposition}
\newtheorem{lemma}[theorem]{Lemma}
\theoremstyle{remark}
\newtheorem{remark}[theorem]{Remark}

\DeclareMathOperator{\Aut}{Aut}
\DeclareMathOperator{\Syl}{Syl}
\DeclareMathOperator{\pr}{pr}
\DeclareMathOperator{\Stab}{Stab}
\newcommand{\gen}[1]{\langle #1\rangle}
\newcommand{\GammaLA}{\Gamma_{L,A}}

\title[Subgroups of coprime index]{Generation of finite groups from subgroups of coprime index}
\author{Richie Sater}
\address{Independent Researcher, United States}
\email{richiesater@gmail.com}
\urladdr{https://orcid.org/0009-0007-9051-8207}
\date{August 11, 2026}
\subjclass[2020]{20D10, 20D30, 20P05}
\keywords{finite group, generators, crown-based power, probabilistic generation, presentation rank}

\begin{document}

\begin{abstract}
Let $d(G)$ denote the least size of a generating set of a finite group $G$.
We prove that if $G$ has a family $\mathcal H$ of subgroups such that
$d(H)\leq d$ for every $H\in\mathcal H$ and
\[
  \gcd\{\lvert G:H\rvert:H\in\mathcal H\}=1,
\]
then $d(G)\leq d+1$.  This gives an affirmative answer to Kourovka
Problem~21.87.  The proof reduces a minimal counterexample to a critical
crown-based power with nonabelian socle.  An exact crown multiplicity formula
and a uniform lower bound for conditional generation give a lower bound for
the number of crown factors.  A subgroup containing a Sylow $2$-subgroup gives
the contradictory upper bound, via a pointwise centralizer estimate for Sylow
$2$-subgroups of finite simple groups.
\end{abstract}

\maketitle

\section{Introduction}

For a finite group $X$, write $d(X)$ for the minimum cardinality of a
generating set of $X$.  Kourovka Problem~21.87, proposed by Lucchini, asks
whether a finite group with $d$-generated subgroups whose indices have no
common divisor must be $(d+1)$-generated \cite[Problem 21.87,
p.~177]{kourovka21}.  Kov\'acs and Sim proved this for soluble groups
\cite[Theorem~2]{KovacsSim1991}; Lucchini subsequently proved the general
$d+2$ bound \cite{Lucchini2000}.  We prove the sharp general bound.

\begin{theorem}\label{thm:main}
Let $d$ be a nonnegative integer, let $G$ be a finite group, and suppose that
$G$ has a nonempty family $\mathcal H$ of subgroups satisfying
\[
 d(H)\leq d\quad(H\in\mathcal H),
 \qquad
 \gcd\{\lvert G:H\rvert:H\in\mathcal H\}=1.
\]
Then $d(G)\leq d+1$.
\end{theorem}

The bound is best possible already when $d=1$.  Indeed, in $S_3$ the cyclic
subgroups $H_2=\gen{(12)}$ and $H_3=\gen{(123)}$ have
$\lvert S_3:H_2\rvert=3$ and $\lvert S_3:H_3\rvert=2$, while $d(S_3)=2$.

The proof combines six published inputs.  Lucchini's augmentation-ideal
bound gives the required module-generator estimate
\cite[Theorem~2 and Proposition~2]{Lucchini1992}.  Roggenkamp's
presentation-rank identity then localizes the only possible obstruction in
positive presentation rank \cite[Theorem~2.1]{Roggenkamp1979}.  A result of
Gruenberg eliminates soluble normal subgroups from a minimal counterexample
\cite[p.~218]{Gruenberg1976}; we use the explicit formulation in
\cite[2.1, p.~209]{Lucchini1990}.  Dalla Volta and Lucchini identify a group
which needs more generators than all its proper quotients as a critical
crown-based power and give its exact multiplicity
\cite[Theorems~1.4 and 2.7]{DallaVoltaLucchini1998}.  Finally, Detomi and
Lucchini prove that the relevant conditional generation probability is at
least $1/2$ \cite{DetomiLucchini2013}.  The new ingredient is to use the
Sylow $2$-subgroup supplied by the index hypothesis to pack distinct
automorphism orbits inside one fiber.  Guralnick's theorem that a finite
nonabelian simple group is generated by a Sylow $2$-subgroup and an
involution \cite{Guralnick1986} supplies exactly the matching factor $1/2$.
The form of Guralnick's theorem used here is also stated explicitly in
\cite[Introduction]{BurnessGuralnick2024}.

The published inputs from Dalla Volta--Lucchini, Detomi--Lucchini, and
Guralnick depend on the classification of finite simple groups.  No
computational classification or finite enumeration is used in the proof.

\section{Preliminaries}

\subsection{The index hypothesis}

\begin{lemma}[Quotient inheritance]\label{lem:quotient}
If $N\unlhd G$, then $G/N$ satisfies the hypothesis of
\cref{thm:main} with the same value of $d$.
\end{lemma}

\begin{proof}
For $H\in\mathcal H$, the subgroup $HN/N$ is generated by at most $d$
elements and
\[
 \lvert G/N:HN/N\rvert=\lvert G:HN\rvert
 \quad\text{divides}\quad \lvert G:H\rvert.
\]
Consequently every common divisor of the image indices divides every
original index, and hence the image indices have greatest common divisor
one.
\end{proof}

\begin{lemma}[Sylow coverage]\label{lem:sylow-coverage}
The index gcd is one if and only if, for every prime $p$ dividing
$\lvert G\rvert$, some $H\in\mathcal H$ contains a Sylow $p$-subgroup of
$G$.
\end{lemma}

\begin{proof}
There is an $H\in\mathcal H$ for which
$p\nmid\lvert G:H\rvert$.  Thus the $p$-part of $\lvert H\rvert$ equals the
$p$-part of $\lvert G\rvert$, so every Sylow $p$-subgroup of $H$ is a Sylow
$p$-subgroup of $G$.  Conversely, if the index gcd were greater than one,
some prime $p\mid\lvert G\rvert$ would divide every index.  No member of
$\mathcal H$ could then contain a Sylow $p$-subgroup of $G$.
\end{proof}

Because all indices divide $\lvert G\rvert$, the gcd of an arbitrary family
of indices is already attained by a finite subfamily.  Thus no issue of an
infinite gcd is hidden in the statement, and the published augmentation-ideal
theorem may be applied to such a finite subfamily.

\begin{lemma}[Sylow transport]\label{lem:sylow-transport}
Let $P\in\Syl_p(X)$.
\begin{enumerate}
\item If $N\unlhd X$, then $P\cap N\in\Syl_p(N)$.
\item If $\pi:X\to Y$ is surjective, then $\pi(P)\in\Syl_p(Y)$.
\end{enumerate}
\end{lemma}

\begin{proof}
The first assertion is the standard normal-subgroup part of the Sylow
theorems.  Applying it to $\ker\pi$ gives
\[
 |\pi(P)|=\frac{|P|}{|P\cap\ker\pi|}
          =\frac{|X|_p}{|\ker\pi|_p}=|Y|_p,
\]
which proves the second assertion.
\end{proof}

\subsection{Augmentation ideals and critical crowns}

Let $I_X$ denote the augmentation ideal of $\mathbb Z X$, and let
$d_X(I_X)$ be its minimum number of generators as a $\mathbb Z X$-module.
If $X$ is generated by $r$ elements, then $I_X$ is generated as a module by
the corresponding $r$ elements of the form $1-x$.  Lucchini's theorem
therefore yields the following consequence.

\begin{proposition}[Lucchini]\label{prop:augmentation}
Under the hypotheses of \cref{thm:main},
\[
 d_G(I_G)\leq d+1.
\]
\end{proposition}

This is \cite[Theorem~2 and Proposition~2]{Lucchini1992}.

\begin{proposition}[Roggenkamp]\label{prop:presentation-rank}
For every finite group $X$,
\[
 d(X)=d_X(I_X)+\pr(X),
\]
where $\pr(X)$ is the presentation rank of $X$.
\end{proposition}

Roggenkamp proves this in \cite[Theorem~2.1]{Roggenkamp1979}.
Lucchini also records the identity, with attribution to Roggenkamp, in
\cite[p.~146]{Lucchini1992}.

For later use, we record the soluble-normal-subgroup consequence of
Gruenberg's presentation-rank theory.

\begin{proposition}[Gruenberg]\label{prop:pr}
If $X$ is finite, $\pr(X)>0$, and $N\unlhd X$ is soluble, then
$d(X)=d(X/N)$.
\end{proposition}

We use the formulation in \cite[2.1, p.~209]{Lucchini1990}, which cites
\cite[p.~218]{Gruenberg1976}.

We next fix the crown notation.  Let $L$ be a finite monolithic group with
socle $A$.  When $A$ is abelian, assume that $A$ is complemented in $L$.
For $k\geq 1$, define the crown-based power
\[
 L_k=\{(\ell_1,\ldots,\ell_k)\in L^k:
          \ell_1A=\cdots=\ell_kA\}.
\]
Set
\[
 \GammaLA=
 \{\alpha\in\Aut(L):\alpha(\ell)A=\ell A\text{ for every }\ell\in L\}.
\]
Thus $\GammaLA$ is the group of automorphisms of $L$ inducing the identity
on $L/A$.

Let $\phi_X(m)$ be the number of ordered generating $m$-tuples of $X$.
For $m\geq d(L)$, put
\[
 P_{L,A}(m)=
 \frac{\phi_L(m)}{|A|^m\phi_{L/A}(m)}.
\]
This is the probability that an $m$-tuple of lifts generates $L$, conditional
on its image generating $L/A$.

We use the following two published results in exactly this form.

\begin{proposition}[Critical crown formula]\label{prop:crown}
Let $m\geq2$, and suppose that $d(X/N)\leq m$ for every nontrivial normal
subgroup $N$ of a finite group $X$, while $d(X)>m$.  Then
$X\cong L_{f_L(m)}$ for a monolithic group $L$ as above.  If its socle $A$
is nonabelian, then
\begin{equation}\label{eq:crown-formula}
 f_L(m)=1+\frac{P_{L,A}(m)|A|^m}{|\GammaLA|}.
\end{equation}
\end{proposition}

This is \cite[Theorems~1.4 and 2.7]{DallaVoltaLucchini1998}, with the
nonabelian branch of the latter normalized as follows.

In the notation of that source, the nonabelian-socle branch is
\[
 f_L(m)=1+\frac{\phi_L(m)}
 {|{\GammaLA}|\,\phi_{L/A}(m)}.
\]
Substitution of the preceding definition of $P_{L,A}(m)$ gives
\eqref{eq:crown-formula}; thus no probabilistic normalization is implicit in
the quoted formula.

\begin{proposition}[Detomi--Lucchini]\label{prop:probability}
If $L$ is a finite monolithic group with socle $A$ and $m\geq d(L)$, then
\begin{equation}\label{eq:half}
 P_{L,A}(m)\geq\frac12.
\end{equation}
\end{proposition}

This is the main inequality of \cite{DetomiLucchini2013}.

The notation in \cref{prop:crown} is chosen so that $f_L(m)$ is the first
crown multiplicity at which more than $m$ generators are needed.  In
particular, $f_L(m)\geq2$, and $L$ is a proper quotient of
$L_{f_L(m)}$.

\section{Pointwise Sylow centralizers}

For a subgroup $Y\leq X$, the notation $C_{\Aut(X)}(Y)$ means the pointwise
centralizer of $Y$ in $\Aut(X)$.

\begin{lemma}[Simple socle factor]\label{lem:simple-centralizer}
Let $S$ be a finite nonabelian simple group and let
$P\in\Syl_2(S)$.  Then
\begin{equation}\label{eq:simple-centralizer}
 |C_{\Aut(S)}(P)|\leq \frac{|S|}{2}.
\end{equation}
\end{lemma}

\begin{proof}
By Guralnick's theorem \cite{Guralnick1986}, there are
$P_0\in\Syl_2(S)$ and an involution $t_0\in S$ such that
$S=\gen{P_0,t_0}$.  Sylow conjugacy gives $s\in S$ with $P_0^s=P$;
after conjugating the generating pair and setting $t=t_0^s$, we have
$S=\gen{P,t}$.  Put $D=C_{\Aut(S)}(P)$.  The $S$-conjugacy class
$t^S$ is $D$-invariant.  Indeed, $t$ is conjugate in $S$ to an involution
$u\in P$; every element of $D$ fixes $u$, so it preserves $u^S=t^S$.

The stabilizer $D_t$ is trivial: an element of $D_t$ fixes $P$ pointwise and
fixes $t$, hence fixes the generating group $\gen{P,t}=S$ pointwise.  The
$D$-orbit of $t$ therefore has size $|D|$ and lies in $t^S$.  Since
$\lvert C_S(t)\rvert\geq2$,
\[
 |D|\leq |t^S|=\lvert S:C_S(t)\rvert\leq |S|/2.
\]
\end{proof}

\begin{lemma}[Monolithic group]\label{lem:monolithic-centralizer}
Let $L$ be monolithic with nonabelian socle $A$, and let
$R\in\Syl_2(L)$.  Then
\begin{equation}\label{eq:monolithic-centralizer}
 |C_{\GammaLA}(R)|\leq |A|/2.
\end{equation}
\end{lemma}

\begin{proof}
Write
$A=S_1\times\cdots\times S_n$, where the $S_i$ are isomorphic finite
nonabelian simple groups.  Restriction gives an embedding
$\GammaLA\hookrightarrow\Aut(A)$.  To see injectivity, suppose that
$\alpha\in\GammaLA$ fixes $A$ pointwise.  For $x\in L$, write
$\alpha(x)=xc$ with $c\in A$.  Comparing $\alpha(xax^{-1})$ with
$\alpha(x)\alpha(a)\alpha(x)^{-1}$ for $a\in A$ shows that
$c\in Z(A)=1$.  Hence $\alpha=1$.

Now
\[
 R\cap A=P_1\times\cdots\times P_n,
 \qquad P_i\in\Syl_2(S_i).
\]
Every $P_i$ is nontrivial.  An automorphism of $A$ which fixes
$R\cap A$ pointwise cannot permute the factors, because a nonidentity
element supported in $P_i$ would be moved to a different support.  By
\cref{lem:simple-centralizer}, restriction to the factors gives
\[
 |C_{\GammaLA}(R)|
 \leq \prod_{i=1}^n |C_{\Aut(S_i)}(P_i)|
 \leq (|S_1|/2)^n
 \leq |A|/2.
\]
\end{proof}

\section{Proof of the main theorem}

\begin{proof}[Proof of \cref{thm:main}]
If $d=0$, every member of $\mathcal H$ is trivial.  The gcd condition then
forces $|G|=1$, so assume $d\geq1$ and put $m=d+1\geq2$.

Suppose, for a contradiction, that the theorem is false.  For this fixed
value of $d$, choose a counterexample $G$ of least order.  By
\cref{lem:quotient}, every proper quotient of $G$ satisfies the same
hypothesis with the same $d$.  Minimality gives
\begin{equation}\label{eq:proper-quotients}
 d(G/N)\leq m
 \qquad(1\neq N\unlhd G),
\end{equation}
whereas $d(G)>m$.

By \cref{prop:augmentation}, $d_G(I_G)\leq m$, so
\cref{prop:presentation-rank} implies $\pr(G)>0$.  If $G$ had a nontrivial
soluble normal subgroup $N$, then \cref{prop:pr} and
\eqref{eq:proper-quotients} would give
$d(G)=d(G/N)\leq m$, a contradiction.  Thus
\begin{equation}\label{eq:no-soluble-normal}
 G\text{ has no nontrivial soluble normal subgroup.}
\end{equation}

Apply \cref{prop:crown}.  We have
\[
 G\cong L_k,
 \qquad k=f_L(m),
\]
where $L$ is monolithic with socle $A$.  The socle $A^k$ of $G$ is a
nontrivial normal subgroup, so \eqref{eq:no-soluble-normal} excludes the
abelian-socle case.  These are the two exhaustive socle branches in
\cref{prop:crown}, so \(A\) is nonabelian.  Also \(L\) is a proper quotient
of \(G\), so \(d(L)\leq m\).  Combining
\cref{prop:crown,prop:probability} gives
\begin{equation}\label{eq:lower-k}
 k
 =1+\frac{P_{L,A}(m)|A|^m}{|\GammaLA|}
 \geq 1+\frac{|A|^m}{2|\GammaLA|}.
\end{equation}

Each simple direct factor of $A$ contains the involution supplied by
Guralnick's theorem, so $A$, and hence $G$, has even order.  By
\cref{lem:sylow-coverage}, there is an $H\in\mathcal H$ containing a Sylow
$2$-subgroup $T$ of $G$.  After padding a generating tuple with identities,
write
\[
 H=\gen{h_1,\ldots,h_d},
 \qquad
 h_i=(h_{i1},\ldots,h_{ik})\in L_k.
\]
For $1\leq j\leq k$, set
\[
 \mathbf h_j=(h_{1j},\ldots,h_{dj})\in L^d.
\]
All these tuples have the same image in $(L/A)^d$.  They therefore belong
to one fiber $\Omega\subseteq L^d$ of cardinality $|A|^d$.  The group
$\GammaLA$ acts componentwise on $\Omega$.

We claim that $\mathbf h_1,\ldots,\mathbf h_k$ belong to distinct
$\GammaLA$-orbits.  Suppose instead that, for distinct $r,s$, there is an
$\alpha\in\GammaLA$ such that
\[
 h_{is}=\alpha(h_{ir})\qquad(1\leq i\leq d).
\]
It follows by evaluating words in the generators that every $h\in H$
satisfies $h_s=\alpha(h_r)$.  But, by
\cref{lem:sylow-transport}, $T\cap A^k$ is a Sylow $2$-subgroup of $A^k$,
hence is a direct product of $k$ nontrivial Sylow $2$-subgroups, one in each
coordinate.  Since
$T\leq H$, the group $H$ contains an element of $T\cap A^k$ supported only
in coordinate $r$.  The displayed relation would send its nonidentity
$r$-coordinate to its identity $s$-coordinate, which is impossible.  This
proves the claim.

Let $T_j$ denote the image of $T$ under the $j$th coordinate projection
$L_k\to L$.  This projection is surjective, so
\cref{lem:sylow-transport} gives $T_j\in\Syl_2(L)$.  The
stabilizer of $\mathbf h_j$ in $\GammaLA$ fixes
$\gen{h_{1j},\ldots,h_{dj}}$ pointwise.  This subgroup contains $T_j$ because
$T\leq H$.  By \cref{lem:monolithic-centralizer},
\[
 |\Stab_{\GammaLA}(\mathbf h_j)|\leq |A|/2.
\]
Thus each of the $k$ distinct orbits has at least
$2|\GammaLA|/|A|$ elements.  Packing them into $\Omega$ gives
\begin{equation}\label{eq:upper-k}
 k\frac{2|\GammaLA|}{|A|}\leq |A|^d,
 \qquad\text{and hence}\qquad
 k\leq \frac{|A|^{d+1}}{2|\GammaLA|}
   =\frac{|A|^m}{2|\GammaLA|}.
\end{equation}
The upper bound \eqref{eq:upper-k} contradicts the strict lower bound
\eqref{eq:lower-k}.  This proves $d(G)\leq d+1$.
\end{proof}

\begin{remark}
The sharp point is the comparison between \eqref{eq:lower-k} and
\eqref{eq:upper-k}.  The conditional probability supplies one factor
$1/2$, while the pointwise Sylow-centralizer estimate supplies the matching
factor in the orbit count.  The additional $1$ in the exact critical-crown
formula rules out equality.
\end{remark}

\section*{Dependencies and disclosure}

The proof is entirely deductive from the published results cited above.  No
exploratory computation is used as evidence.  The cited results on critical
crowns, conditional generation, and generation of simple groups depend on
the classification of finite simple groups.  A dated literature search found
no earlier paper asserting this conclusion, but that search is not a claim
of priority.  This manuscript has not undergone external specialist review.

The author used OpenAI Codex for literature organization, proof stress-testing,
and editorial assistance.  Responsibility for the mathematical claims and
any submitted version remains with the author.

\bibliographystyle{amsplain}
\bibliography{references}

\end{document}
